% Source: Rsquared_coeficient.tex
\begin{notechapter}
\chapter{The Coefficient of Determination $R^2$}\label{ch:r-squared}

The coefficient of determination, denoted as \( R^2 \), is a statistical 
measure that assesses the goodness of fit of a regression model. It 
indicates the proportion of the variance in the dependent variable that 
is predictable from the independent variable(s).

\section{Definition}
\( R^2 \) is defined as the ratio of the explained variance to the total
variance. For a least-squares linear model with an intercept, evaluated on
the data it was fitted on, it ranges from 0 to 1, where:
\begin{itemize}
    \item \( R^2 = 0 \) indicates that the model explains none of the 
          variability of the response data around its mean.
    \item \( R^2 = 1 \) indicates that the model explains all the 
          variability of the response data around its mean.
\end{itemize}

\section{Formula}
For any regression model, \( R^2 \) is calculated as:
\[
R^2 = 1 - \frac{SS_{\text{res}}}{SS_{\text{tot}}}
\]
where:
\begin{itemize}
    \item \( SS_{\text{res}} \) is the sum of squares of residuals:
          \[
          SS_{\text{res}} = \sum_{i=1}^n (y_i - \hat{y}_i)^2
          \]
    \item \( SS_{\text{tot}} \) is the total sum of squares:
          \[
          SS_{\text{tot}} = \sum_{i=1}^n (y_i - \bar{y})^2
          \]
\end{itemize}
Here:
\begin{itemize}
    \item \( y_i \) are the observed values.
    \item \( \hat{y}_i \) are the predicted values.
    \item \( \bar{y} \) is the mean of the observed values.
        \item \( n \) is the number of observations.
\end{itemize}
For ordinary least squares with an intercept, the total sum of squares
decomposes as \( SS_{\text{tot}} = SS_{\text{reg}} + SS_{\text{res}} \) with
\( SS_{\text{reg}} = \sum_{i=1}^n (\hat{y}_i - \bar{y})^2 \), so that
\( R^2 = SS_{\text{reg}}/SS_{\text{tot}} \in [0,1] \): this is the ``explained
variance over total variance'' reading of the definition. In this case
\( R^2 \) is also the squared sample correlation between \( y \) and
\( \hat{y} \). Outside this setting (no intercept, non-linear model, or
evaluation on held-out data) the decomposition fails and \( R^2 \) can be
negative: \( R^2 < 0 \) means that the model predicts worse than the
constant predictor \( \bar{y} \).

\section{Interpretation}
\begin{itemize}
    \item \textbf{High \( R^2 \)}: A high \( R^2 \) value (close to 1) 
          suggests that the model explains a large portion of the 
          variance in the dependent variable, indicating a good fit.
    \item \textbf{Low \( R^2 \)}: A low \( R^2 \) value (close to 0) 
          suggests that the model does not explain much of the variance 
          in the dependent variable, indicating a poor fit.
\end{itemize}

\section{Limitations}
\begin{itemize}
        \item \textbf{Overfitting}: for nested least-squares models, adding a
          predictor can never decrease the training \( R^2 \), so
          \( R^2 \) can be misleadingly high in models with many
          predictors, even if those predictors are not meaningful.
    \item \textbf{Model dependence}: \( R^2 \) only measures how well the
          \emph{fitted} model accounts for the variance. A low \( R^2 \)
          for a linear model does not mean that \( X \) and \( Y \) are
          unrelated (the relation may be non-linear), and a high \( R^2 \)
          does not prove that the model is correctly specified.
    \item \textbf{Outliers}: \( R^2 \) can be sensitive to outliers, 
          which can disproportionately affect the sum of squares.
\end{itemize}

\section{Adjusted \( R^2 \)}
To account for the number of predictors in the model, the adjusted \( R^2 \) 
is often used. It adjusts the \( R^2 \) value based on the number of 
predictors and the sample size:
\[
\text{Adjusted } R^2 = 1 - \left( \frac{(1 - R^2)(n - 1)}{n - p - 1} \right)
\]
where:
\begin{itemize}
    \item \( n \) is the sample size.
    \item \( p \) is the number of predictors.
\end{itemize}

\section{Use Cases}
\begin{itemize}
    \item \textbf{Model Comparison}: \( R^2 \) is useful for comparing 
          the explanatory power of different models.
    \item \textbf{Performance Metric}: It serves as a performance metric 
          for regression models, though it should be used in conjunction 
          with other metrics like RMSE or MAE.
\end{itemize}

\section{Conclusion}
While \( R^2 \) is a widely used metric to evaluate the performance of 
regression models, it should be interpreted with caution, especially in 
the context of model complexity and the nature of the data. Adjusted \( R^2 \) 
and other metrics can provide a more comprehensive understanding of 
model performance.
\end{notechapter}
